\section{大动量展开对短距离展开}

文献 \cite{Ji:2013dva} 引入的式 (5) 中的欧氏关联子
也在坐标空间因子化（CSF）~\cite{Radyushkin:2017cyf} 中被考虑；CSF 最早见于用流-流关联子研究介子 DA 的早期工作
\cite{Braun:2007wv}（另见~\cite{Ma:2017pxb}）。
该关联子可以按展开参数 $(z\Lambda_{\rm QCD})^2$ 用 LF 关联来因子化。这一形式体系
天然适合计算 PDF 的矩或短距离 LF 关联。但要获得完整的部分子物理，还必须同时考虑外部动量的约束 $P^z\sim 1/z \gg \Lambda_{\rm QCD}$。
这与文献 \cite{Ji:2014gla} 的观察一致：必须使用大动量才能覆盖 PDF 的完整动力学范围，这需要 $\lambda$ 中长程关联的信息。
尽管二者在形式上等价~\cite{Ji:2017rah,Izubuchi:2018srq,Radyushkin:2017lvu}，
某些解析匹配计算在坐标空间中进行或许更为方便。不出所料，CSF 分析要得到 PDF 同样需要相同的 LaMET 格点数据。
名义上 CSF 也可纳入小 $P^z$ 数据，但这些信息其实已包含在更小 $z$ 处的大 $P^z$ 数据中。

CSF 展开以欧氏距离 $z$ 表述，要求 $z$ 很小，即
\begin{equation}
         z\ll 1/\Lambda_{\rm QCD}\ ,
\end{equation}
以保证微扰论有效和领头扭度主导。设领头扭度近似适用的最大 $z$ 为 $z_{\rm LT}$（比如高扭度贡献达到 $\sim20\%$ 水平时对应的 $z$ 值），那么在展开之前先行减除可能的线性发散时，小展开参数为 $(z_{\rm LT}\Lambda_{\rm QCD})^2\ll 1 $。因此只有 $[0, z_{\rm LT}]$ 范围内 $O(z)$ 的矩阵元才具有以领头扭度部分子物理表述的简单诠释。

于是一个有趣的问题是：$z_{\rm LT}$ 取什么值？
若取 $\Lambda_{\rm QCD}\sim 300$ MeV，并把 $z_{\rm LT}\Lambda_{\rm QCD}=1/2\sim 1/3$ 视为小参数，则估计 $z_{\rm LT}$ 约为 $0.25 \sim 0.33$ fm。上限大概是 0.4 fm。把矩阵元 $\langle P=0| O(z)|P=0\rangle$ 或 $\langle \Omega| O(z)|\Omega\rangle$——两者都被提议用于重正化裸准 LF 关联~\cite{Radyushkin:2017cyf,Braun:2018brg,Li:2020xml}——与其 OPE 中的领头扭度贡献作比较，可以较好地估计 $z_{\rm LT}$。

我们以同位旋矢情形的零动量矩阵元为例。在 $\MS$ 方案中它具有如下形式的短距离展开~\cite{Radyushkin:2017lvu,Izubuchi:2018srq}
\begin{align}
			\tilde{h}(z,\mu,P^z\!=\!0)&= {1\over 2M}\langle P=0| \bar{\psi}(z)\gamma^0 W(z,0)\psi(0)|P=0\rangle \nn\\
			& = c_0(\mu^2 z^2) a_0 + {\cal O}(z^2\Lambda_{\rm QCD}^2)\,,
\end{align}
其中 $W(z,0)$ 是类空直线规范连接，$\mu$ 是 $\MS$ 重正化标度，$a_0=1$ 是相应 twist-2 PDF 守恒的最低矩。单圈 Wilson 系数 $c_0=1/  Z^{\overline{\rm MS}}_{\rm ratio}$，其中 $Z^{\overline{\rm MS}}_{\rm ratio}$ 由 \eq{zmsr} 给出~\cite{Izubuchi:2018srq}，
两圈结果见 ~\cite{Li:2020xml}。

根据 \eq{ren}，质量减除后的矩阵元所含的对数发散与 $z$ 无关，不应构成格点微扰论中的显著修正。因此可以粗略地将其 OPE 近似为把 $\mu$ 替换成 $1/a$：
\begin{align}\label{eq:subtracted}
	e^{-\delta m |z|}\tilde{f}(z,a,P^z\!=\!0) \!=\! c_0(z^2/a^2) \!+\! {\cal O}(z^2\Lambda_{\rm QCD}^2, a^2\!/\!z^2)\,,
\end{align}
其中格点离散化效应预期为 ${\cal O}(a^2/z^2)$。
由于格点矩阵元在 $z\to0$ 时收敛——这与 $\MS$ OPE 的对数发散行为相反——我们预期上述近似在离散化与高扭度效应都被压低的范围 $a < |z| < z_{\rm LT}$ 内可靠。注意，格点 OPE 通常因 Lorentz 对称性破缺与算符混合而变得复杂。不过，既然 $P^z=0$ 时唯一的领头扭度贡献来自守恒矢量流，这里可以忽略这些效应。

\fig{ope} 中绘出了质量重正化的 π 介子格点矩阵元。裸的格点矩阵元来自近期在一个系综（$a=0.06$ fm、π 介子质量 $m_\pi=300$ MeV）上对 π 介子价 PDF 的计算~\cite{Gao:2020ito}。在同一格点系综上，Wilson 线质量修正 $\delta m$ 由夸克-反夸克势拟合得到~\cite{Izubuchi:2019lyk}，其值以格点单位表示为 $a \delta m = -0.1568$。图中领头扭度贡献分别画出了次领头阶（NLO）修正以及固定 $\alpha_s$ 下 NLO 修正加重求和的情形：
\begin{align}
	\left[ 1 + {\alpha_s(1/a)C_F\over 2\pi}  {5\over2}\right]\left({z^2\over 4e^{-2\gamma_E}a^2}\right)^{{\alpha_s(1/a)C_F\over 2\pi}{3\over 2}}\,.
\end{align}
这里选取标度 $1/a$ 处 $\MS$ 方案的 $\alpha_s$ 作为 OPE 输入，这应比裸格点耦合有更好的收敛性~\cite{Lepage:1992xa}。为估计 $\alpha_s$ 选择带来的不确定性，我们把 $\MS$ 标度从 $1/(2a)$ 变到 $2/a$。虽然定义格点上的 $\alpha_s$ 还须执行标准的改进程序，但我们预期这不会改变下文的结论。

可以看到，当 $z\le a$ 时，由于离散化效应，格点结果与领头扭度近似显著不同。随 $z$ 增大，二者吻合变好。然而当 $z\ge 0.3$ fm 时，格点结果开始明显偏离领头扭度近似，表明高扭度贡献变得显著。因此我们可以粗略估计 $z_{\rm LT}\sim 0.3$ fm。

也可以考察研究得更成熟的重夸克势情形。众所周知，静态重夸克势同时接受微扰与非微扰贡献。微扰静势已知到 ${\rm N^3LO}$ 阶~\cite{Smirnov:2009fh}
并可用 QCD 跑动耦合常数表达。文献 ~\cite{Bazavov:2012ka,Bazavov:2014soa,Bazavov:2019qoo} 从短距离静态能的格点计算中提取了跑动耦合常数。${\rm N^3LL}$ 微扰结果在 $r \sim 0.2$ ${\rm fm}$ 以内与格点数据吻合良好。然而众所周知，静势的微扰级数存在 renormalon 模糊~\cite{Beneke:1998rk,Hoang:1998nz}，在大距离处失效。非微扰重夸克势可用格点 QCD 模拟，且众所周知在大距离由形如 $\sigma r$ 的线性项主导。唯象上，静势可由 linear+Coulumb 型 QCD 静势很好地近似：$V(r)=-e/r+\sigma r$，其中 $e= 0.25\sim 0.5$ 而 $\sqrt{\sigma} \approx 477$ ${\rm MeV}$~\cite{Eichten:1974af,Bali:2000gf,Aubin:2004wf}。当微扰论即将失效时，微扰贡献 $-e/r$ 与禁闭贡献 $\sigma r$ 应为同一量级，由此确定 $r_c \approx\sqrt{e/\sigma}$ 约为 $0.2\sim 0.3$ ${\rm fm}$，与 ~\cite{Bali:1999ai,Necco:2001gh,Bazavov:2012ka} 的结果一致。边界 $z_{\rm LT}$ 应为同一量级。

\begin{figure}[t]
\centering
\includegraphics[width=1\columnwidth]{images/ZLT_pion}
\caption{质量重正化的 π 介子格点矩阵元~\cite{Gao:2020ito} 与其领头扭度近似（分别为 NLO 及固定 $\alpha_s$ 下 NLO$+$LL 修正）的比较。
强耦合为 $\alpha_s(1/a)=0.242$，误差带由把 $\alpha_s$ 从 $\alpha_s(1/(2a)$ 变到 $\alpha_s(2/a)$ 得到。这表明领头扭度近似在 $z_{\rm LT}\sim 0.3$ fm 处变得不可靠。}
\label{fig:ope}
\end{figure}

\begin{figure}[t]
\includegraphics[width=1\columnwidth]{images/potential}
\caption{$\alpha_s^3$ 阶微扰势与 linear+Coulumb 型 QCD 静势的比较，取自~\cite{Bali:2000gf,Aubin:2004wf}。阴影区域对应重正化标度 $\mu$ 从 $r/2$ 到 $3r/2$ 的微扰预言，黑线对应 $\mu=1/r$。由图清楚可见，超出 $0.2$ 至 $0.3$ ${\rm fm}$ 之后，微扰势开始偏离完整的非微扰结果并变得不可靠。}
\label{fig:poten}
\end{figure}

利用上面估计的 $z_{\rm LT}$，还可以定义
\begin{equation}
       \lambda_{\rm LT} = P^zz_{\rm LT},
\end{equation}
于是准 LF 距离 $[0, \lambda_{\rm LT}]$ 内的矩阵元可以用 \eq{coordfact} 的匹配公式来提取部分子分布~\cite{Radyushkin:2017cyf,Ji:2017rah,Izubuchi:2018srq}。
因此，坐标空间途径适用于在有限范围（LF 距离介于 $[0,\lambda_{\rm LT}]$）内提取 LF 关联函数。

CSF 方法也被用于夸克双线型构成的流乘积~\cite{Detmold:2005gg,Braun:2007wv,Bali:2017gfr,Bali:2018spj,Ma:2017pxb,Detmold:2018kwu,Sufian:2019bol,Sufian:2020vzb}。
对于带 Wilson 线的夸克与胶子双线型，CSF 展开中幂次发散的重正化更易处理。特别地，用零动量矩阵元作分母的一版比值方法可以消除格点矩阵元中的幂次发散~\cite{Radyushkin:2017cyf,Orginos:2017kos}。另一方面，流乘积经傅里叶变换到动量空间后也可用于 LaMET 展开~\cite{Ji:2020ect}。

然而，CSF 不允许直接计算 $x$ 分布，因为那需要所有 LF 距离处的 LF 关联。$z\le z_{\rm LT}$ 的要求使得在当代格点资源的最大动量下，傅里叶变换所需的大 $\lambda$ 值难以企及。于是要重构 PDF，只能像唯象拟合中那样参数化 $x$ 依赖的函数形式，然后把它逆傅里叶变换回坐标空间，并对有限范围的
LF 关联进行拟合~\cite{Orginos:2017kos,Karpie:2018zaz,Joo:2019bzr,Joo:2019jct,Sufian:2019bol,Joo:2020spy,Sufian:2020vzb,Bhat:2020ktg,Gao:2020ito,Fan:2020cpa}。
这一过程几乎不受控制，因为很难估计来自 PDF 参数化或其假设的系统不确定性。正如迄今文献中的拟合所显示的那样~\cite{Joo:2019bzr,Joo:2019jct,Sufian:2019bol,Joo:2020spy,Gao:2020ito,Fan:2020cpa}：要么端点区域的误差越变越小，要么在某些中等 $x$ 值处误差收缩到几乎为零。这意味着存在未计入的、来自所用特定模型伪迹的系统误差；这也体现在它们与采用类似参数化的全局拟合互不相符上。

换一个角度看，上述做法相当于对 LF 关联短程行为与长程行为之间的某种关联作出假设（或建模）。这种假设没有第一性原理的基础；而且完全可能出现这样的情形：有限 $\lambda$ 范围内的格点数据可以被不止一个参数化同等良好地拟合，而这些参数化的渐近行为截然不同~\cite{Sufian:2020vzb}。

尽管难于对 $x$ 分布给出受控的计算，CSF 方法允许借助 OPE 以模型无关的方式提取 Mellin 矩。不过，在有限的 $\lambda$ 范围内，LF 关联将只对最低几阶矩敏感；这与 $x$ 空间的预言能力相关，但并无直接的一一对应关系。

为更直接地比较动量空间与坐标空间两种方法，考虑如下例子。设定义准 PDF 的准 LF 关联的行为为
\begin{align}
\tilde h(\lambda,z)=h(\lambda)e^{-m|z|} \ ,
\end{align}
其中 $h(\lambda)$ 是光锥关联子。指数因子 $e^{-m|z|}$（$m\sim \Lambda_{\rm QCD}$）用于模拟高扭度贡献。由该式清楚可见：若停留在位置空间，CSF 只有在 $z\ll 1/m$ 时才准确。可用的 $\lambda$ 范围因而远小于 $P^z/m$，这意味着能够触及的矩数远少于 $P^z/m$。

从上述讨论可以清楚地看到，动量空间与坐标空间展开是不同的展开方案。即便二者在无穷动量极限下等价，它们在有限动量 $P^z$ 下并不相同。
后者直接在坐标空间过滤信息：人们得到的是 LF 距离有限范围内的部分子关联，对应的矩数由 $1/P^z$ 控制。前者则利用全部坐标空间信息，在动量空间通过 $1/(y^2(P^z)^2)\ll 1$ 与 $1/((1-y)^2(P^z)^2)\ll 1$ 过滤高扭度物理。
因此在系统控制误差之下，人们得到的是 $x$ 某个区间内的部分子分布，可直接与实验数据比较。

最后我们指出，准 PDF 微扰修正的相对大小依赖于 $x$：通常在小 $x$ 与大 $x$ 区域观察到更大的修正。另一方面，尽管坐标空间中微扰修正的大小对有限 $\lambda$ 通常不大，它仍会在动量空间的端点区域导致显著的修正。